% from ega1-10-fr.tex
\oldpage[I]{212}
\begin{corollary}[10.14.7]
\label{I.10.14.7-fr}
Soient $X$ un préschéma usuel localement noethérien, $Y$ un sous-préschéma
fermé de $X$, $j$ l'injection canonique $Y\to X$, $X'$ une partie fermée de
$X$ et $Y'=Y\cap X'$~; alors
$\widehat{j}:Y_{/Y'}\to X_{/X'}$ est une immersion fermée, et pour tout
$\mathscr{O}_Y$-Module cohérent $\mathscr{F}$, on a
\[
  \widehat{j}_*(\mathscr{F}_{/Y'})=(j_*(\mathscr{F}))_{/X'}.
\]
\end{corollary}

Comme $Y'=j^{-1}(X')$, il suffit d'utiliser
(\hyperref[I.10.9.9-fr]{10.9.9}) et d'appliquer
(\hyperref[I.10.14.5-fr]{10.14.5}) et
(\hyperref[I.10.14.6-fr]{10.14.6}).

\subsection{Préschémas formels séparés.}
\label{subsection:I.10.15-fr}

\begin{definition}[10.15.1]
\label{I.10.15.1-fr}
Soient $\mathfrak{S}$ un préschéma formel, $\mathfrak{X}$ un
$\mathfrak{S}$-préschéma formel,
$f:\mathfrak{X}\to\mathfrak{S}$ le morphisme structural. On appelle
morphisme diagonal
$\Delta_{\mathfrak{X}|\mathfrak{S}}:\mathfrak{X}\to
\mathfrak{X}\times_{\mathfrak{S}}\mathfrak{X}$ (noté aussi
$\Delta_{\mathfrak{X}}$) le morphisme
$(1_{\mathfrak{X}},1_{\mathfrak{X}})_{\mathfrak{S}}$. On dit que
$\mathfrak{X}$ est séparé sur $\mathfrak{S}$, ou un
$\mathfrak{S}$-schéma formel, ou que $f$ est un morphisme séparé, si l'image
par $\Delta_{\mathfrak{X}}$ de l'espace sous-jacent $\mathfrak{X}$ est une
partie fermée de l'espace sous-jacent à
$\mathfrak{X}\times_{\mathfrak{S}}\mathfrak{X}$. On dit qu'un préschéma
formel $\mathfrak{X}$ est séparé, ou est un schéma formel, s'il est séparé
sur $\mathbf{Z}$.
\end{definition}

\begin{proposition}[10.15.2]
\label{I.10.15.2-fr}
Supposons que les préschémas formels $\mathfrak{S}$, $\mathfrak{X}$ soient
respectivement limites inductives des suites $(S_n)$, $(X_n)$ de
préschémas usuels, et que le morphisme
$f:\mathfrak{X}\to\mathfrak{S}$ soit limite inductive d'une suite de
morphismes $f_n:X_n\to S_n$. Pour que $f$ soit séparé, il faut et il suffit
que le morphisme $f_0:X_0\to S_0$ le soit.
\end{proposition}

En effet, $\Delta_{\mathfrak{X}|\mathfrak{S}}$ est alors limite inductive de
la suite de morphismes $\Delta_{X_n|S_n}$
(\hyperref[I.10.7.4-fr]{10.7.4}), et l'image par
$\Delta_{\mathfrak{X}|\mathfrak{S}}$ de l'espace sous-jacent
$\mathfrak{X}$ (resp. l'espace sous-jacent
$\mathfrak{X}\times_{\mathfrak{S}}\mathfrak{X}$) est identique à l'image
par $\Delta_{X_0|S_0}$ de l'espace sous-jacent $X_0$ (resp. à l'espace
sous-jacent $X_0\times_{S_0}X_0$)~; d'où la conclusion.

\begin{proposition}[10.15.3]
\label{I.10.15.3-fr}
On suppose dans ce qui suit que tous les préschémas formels (resp.
morphismes de préschémas formels) considérés soient limites inductives de
suites de préschémas usuels (resp. de morphismes de préschémas usuels).
\begin{enumerate}
  \item[\rm{(i)}] Le composé de deux morphismes séparés est séparé.
  \item[\rm{(ii)}] Si $f:\mathfrak{X}\to\mathfrak{X}'$,
  $g:\mathfrak{Y}\to\mathfrak{Y}'$ sont deux $\mathfrak{S}$-morphismes
  séparés, $f\times_{\mathfrak{S}}g$ est séparé.
  \item[\rm{(iii)}] Si $f:\mathfrak{X}\to\mathfrak{Y}$ est un
  $\mathfrak{S}$-morphisme séparé, le $\mathfrak{S}'$-morphisme
  $f_{(\mathfrak{S}')}$ est séparé pour toute extension
  $\mathfrak{S}'\to\mathfrak{S}$ du préschéma formel de base.
  \item[\rm{(iv)}] Si le composé $g\circ f$ de deux morphismes est séparé,
  $f$ est séparé.
\end{enumerate}
\emph{(On sous-entend dans cet énoncé que si un même préschéma formel
$\mathfrak{Z}$ intervient plusieurs fois dans une même proposition, on le
considère comme limite inductive de la \emph{même} suite $(Z_n)$ de
préschémas usuels partout où il figure, et les morphismes de $\mathfrak{Z}$
dans un préschéma formel (resp. d'un préschéma formel dans $\mathfrak{Z}$)
comme limites inductives de morphismes des $Z_n$ dans des préschémas usuels
(resp. de préschémas usuels dans les $Z_n$)).}
\end{proposition}

Avec les notations de (\hyperref[I.10.15.2-fr]{10.15.2}), on a en effet
$(g\circ f)_0=g_0\circ f_0$, et
$(f\times_{\mathfrak{S}}g)_0=f_0\times_{S_0}g_0$~; les assertions de
(\hyperref[I.10.15.3-fr]{10.15.3}) sont alors conséquences immédiates de
(\hyperref[I.10.15.2-fr]{10.15.2}) et des assertions correspondantes de
(\hyperref[I.5.5.1-fr]{5.5.1}) pour les préschémas usuels.

Nous laissons au lecteur le soin d'énoncer pour la même espèce de
préschémas formels et de morphismes que dans
(\hyperref[I.10.15.3-fr]{10.15.3}) les propositions correspondant à
(\hyperref[I.5.5.5-fr]{5.5.5}),

